\documentclass[11pt]{article}
\usepackage[utf8]{inputenc}
\usepackage[T1]{fontenc}
\usepackage[italian]{babel}
\usepackage{geometry}
\geometry{a4paper, margin=1in}
\usepackage{xcolor}
\usepackage{listings}
\usepackage{hyperref}

\usepackage{tcolorbox}
\tcbuselibrary{breakable}
\begin{document}

\begin{tcolorbox}[colback=blue!5!white,colframe=blue!75!black,title=\textbf{DOMANDA}, breakable]
2) Modificare i sorgenti per mettere il server che riceve sulla porta 10000 e il client che trasmette dalla propria porta 30000 (ogni modifica dei sorgenti richiede una loro ricompilazione)
\end{tcolorbox}

\begin{tcolorbox}[colback=green!5!white,colframe=green!75!black,title=\textbf{RISPOSTA}, breakable]
\begin{itemize}
    \item Modificato \texttt{serverUDP.c} per metterlo in ascolto sulla porta 10000 (\texttt{\#define SERVERPORT 10000}).
    \item Modificato \texttt{clientUDP.c} per fare in modo che la sua socket venga creata sulla porta 30000 (\texttt{createUDPInterface(30000)}).
    \item Modificato \texttt{clientUDP.c} affinché la \texttt{UDPSend} invii il pacchetto alla porta di destinazione 10000.
\end{itemize}
\end{tcolorbox}

\end{document}
